\documentclass[11pt]{article}

\usepackage[margin=1in]{geometry}
\usepackage{graphicx} % Required for inserting images
\graphicspath{{../manybodyscars/figures/}}
\usepackage{hyperref}
\usepackage{mathtools,amssymb}
\usepackage{amsthm}
\usepackage{mathrsfs}
\usepackage{fancyhdr}
\usepackage{subcaption}
\captionsetup{font=small,labelfont=bf}
\usepackage{float}
\usepackage{mathdots}
\usepackage{physics}
\usepackage{indentfirst}
% \usepackage{fontspec} % Enable with XeLaTeX or LuaLaTeX for custom fonts.

\hypersetup{
  colorlinks=true,
  linkcolor=blue,
  citecolor=blue,
  urlcolor=blue
}

\pagestyle{fancy}
\fancyhf{}
\fancyhead[L]{From Quantum Many-Body Scars to Slow Dynamics}
\fancyhead[R]{Research note}
\fancyfoot[C]{\thepage}
\setlength{\headheight}{14pt}

\newcommand{\Hpxp}{H_{\mathrm{PXP}}}
\newcommand{\Ztwo}{\mathbb{Z}_2}
\newcommand{\Xtilde}{\widetilde{X}}

\title{From Quantum Many-Body Scars to Slow Dynamics\\
  \large Exact Diagonalization of a Deformed PXP Chain}
\author{Zihao Xu}
\date{\today}

\begin{document}

\maketitle

\begin{abstract}
Quantum many-body scars and quantum Griffiths regimes provide two distinct routes
to slow or incomplete thermalization. Scars are atypical eigenstates embedded in an
otherwise thermal spectrum of a clean constrained model, whereas conventional
Griffiths behavior is attributed to rare spatial regions in disordered systems. We
investigate whether these mechanisms can be connected in a clean deformation of the
PXP chain that moves strong $\Ztwo$ spectral weight toward low energy. Exact
diagonalization identifies low-energy eigenstates with anomalously large $\Ztwo$
overlap. Symmetry-resolved spectra determine the canonical energy and heat capacity,
while a thermal-quench calculation yields a connected local $X$-autocorrelation whose
decay is consistent with a power law over the selected time window. These finite-size
results are compatible with slow Griffiths-like dynamics but do not establish a
Griffiths phase. We therefore formulate further tests based on competing dynamical fits,
low-temperature heat-capacity scaling, separation between average and typical
observables, gap distributions, and persistence over a finite parameter window.
\end{abstract}

\section{Introduction}

After exact symmetries and conserved quantities have been resolved, generic
interacting quantum systems are expected to thermalize and exhibit quantum-chaotic
spectral statistics. Their finite-energy-density eigenstates obey the eigenstate
thermalization hypothesis and typically carry volume-law entanglement. Two important
exceptions motivate this work. In strongly disordered spin chains and superconducting
systems, a quantum Griffiths regime can develop when rare spatial regions generate a
broad distribution of relaxation times and exceptionally slow dynamics. In clean
Rydberg-blockaded systems such as the PXP chain, quantum many-body scars instead form
a sparse set of eigenstates with anomalously low entanglement and strong violations of
eigenstate thermalization relative to nearby states.

Both phenomena can retain memory of special initial conditions and produce slow or
reviving dynamics, but their usual microscopic explanations are different. Griffiths
behavior relies on spatial disorder and rare regions, whereas scarred eigenstates are
embedded in the finite-energy-density spectrum of a clean constrained Hamiltonian.
Here we investigate a possible connection between them. A deterministic staggered
deformation of the PXP model shifts the state with the strongest $\Ztwo$ overlap into
the low-energy spectrum, where it can influence low-temperature and long-time
observables. We combine exact diagonalization, canonical thermodynamics, and a local
autocorrelation calculation to test for Griffiths-compatible behavior without
disorder. The present finite-size results identify a candidate route rather than a
complete demonstration of a disorder-free Griffiths phase.

\section{Constrained spin model}

Consider a periodic chain of $L$ spin-$1/2$ degrees of freedom. Let $X_i$ and
$Z_i$ denote Pauli operators, and define
\begin{equation}
  n_i=\frac{1+Z_i}{2},
  \qquad
  P_i=1-n_i=\frac{1-Z_i}{2}.
\end{equation}
The Rydberg blockade restricts the Hilbert space to configurations satisfying
\begin{equation}
  n_i n_{i+1}=0
  \quad \text{for every }i,
\end{equation}
with site labels understood modulo $L$. Within this constrained space, a local spin
flip is equivalent to
\begin{equation}
  \Xtilde_i=P_{i-1}X_iP_{i+1},
\end{equation}
and the standard PXP Hamiltonian is
\begin{equation}
  \Hpxp=\sum_{i=0}^{L-1}\Xtilde_i.
\end{equation}

The model studied here adds staggered diagonal perturbations,
\begin{equation}
  H(g,r)=\Hpxp
  +g\sum_{i=0}^{L-1}(-1)^i
  \left(Z_i+r Z_{i-1}Z_iZ_{i+1}\right).
  \label{eq:hamiltonian}
\end{equation}
The parameter $g$ sets the overall strength of the deformation and $r$ fixes the
relative weight of the three-spin term. The staggering reduces one-site translation
symmetry to translation by two sites. Equation~\eqref{eq:hamiltonian} is implemented
in a custom constrained basis and diagonalized with QuSpin. The deformation is
deterministic and introduces no random couplings. Its role is to tune the spectral
position and stability of the scarred subspace while preserving a clean setting.

The reference product state is
\begin{equation}
  \ket{\Ztwo}=\ket{1010\cdots},
\end{equation}
which is invariant under translation by two sites and therefore belongs to the
zero-momentum sector of that symmetry. Eigenstates with large $\Ztwo$ overlap define
the scarred subspace used below. Bringing this spectral weight toward the bottom of
the spectrum makes the same subspace accessible to low-temperature diagnostics.

\section{Exact diagonalization and scarred eigenstates}

For the spectral calculation, the Hamiltonian is represented in the two-site
translation sector $k=0$. The current calculation uses $L=20$ and obtains the full
dense eigendecomposition
\begin{equation}
  H(g,r)\ket{E_n}=E_n\ket{E_n}.
\end{equation}
For each eigenstate, the principal scar diagnostic is its overlap with the ordered
state,
\begin{equation}
  w_n=\abs{\braket{\Ztwo}{E_n}}^2.
  \label{eq:overlap}
\end{equation}
The spectral location of the initial state can also be summarized by
\begin{equation}
  E_{\Ztwo}=\mel{\Ztwo}{H}{\Ztwo}
  =\sum_n E_n w_n.
\end{equation}

The numerical spectrum contains a small set of eigenstates whose values of $w_n$ are
anomalously large compared with nearby eigenstates. These outliers are quantum
many-body scar candidates. For the deformation of principal interest, $g=-0.4$ and
$r=0.2$, the eigenstate with the largest $\Ztwo$ overlap lies in the low-energy
region of the finite-size spectrum. This is the first ingredient in the proposed 
connection: a selected nonthermal sector is placed where it can affect
low-temperature dynamics. It is not yet a real-space rare region. A central question
is whether the scarred sector acts instead as a rare structure in Hilbert or phase
space.

The half-chain entanglement entropy of every eigenstate is evaluated after lifting the
states from the symmetry-reduced basis to the full constrained basis. It provides an
independent diagnostic of whether large-overlap states are atypical relative to their
spectral neighbors. Unfolded nearest-neighbor level spacings are computed from the
middle two thirds of the spectrum, where edge effects are weaker. Together, the
overlap, entanglement, and level statistics distinguish a sparse scarred subspace from
a global loss of ergodicity or an accidentally integrable spectrum. A systematic
finite-size study is still required for a thermodynamic-limit conclusion.

Figure~\ref{fig:spectrum} collects the spectral diagnostics and the corresponding
quench dynamics. The upper-left panel displays the anomalously large $\Ztwo$
overlaps and places the strongest-overlap state toward the low-energy side of the
spectrum. The lower panels show that the same initial state produces long-lived,
damped oscillations in a local observable and in the return probability. This
spectral selectivity and dynamical memory motivate the thermodynamic and correlation
tests that follow.

\begin{figure}[htbp]
  \centering
  \includegraphics[width=0.96\linewidth]
  {pxp_spectrum_L=20_g=-0.4_r=0.2.png}
  \caption{Spectral and dynamical diagnostics for $L=20$, $g=-0.4$, and
  $r=0.2$. Upper left: overlap $\abs{\braket{\Ztwo}{E_n}}^2$ versus energy;
  the color scale gives the half-chain entanglement entropy, and the dashed line
  marks $E_{\Ztwo}$. Upper right: distribution of unfolded nearest-neighbor level
  spacings. Lower panels: the site-averaged nearest-neighbor correlation
  $L^{-1}\sum_i\expval{Z_iZ_{i+1}}$ and the return probability
  $\abs{\braket{\Ztwo}{\psi(t)}}^2$ after a quench from $\ket{\Ztwo}$.}
  \label{fig:spectrum}
\end{figure}

\section{Canonical energy and heat capacity}

The thermodynamic calculation uses the eigenvalues from all two-site momentum
sectors, rather than only the $k=0$ sector used for the scar plot. If
$E_{n,k}$ denotes the spectrum in an independent momentum block, the partition
function is reconstructed as
\begin{equation}
  Z(\beta)=\sum_{k\in\mathcal{K}_{\mathrm{ind}}}m_k
  \sum_n e^{-\beta E_{n,k}},
  \label{eq:partition}
\end{equation}
where $m_k=2$ for a conjugate pair $k$ and $-k$, while self-conjugate sectors have
$m_k=1$. In practice, the global ground-state energy is subtracted inside the
Boltzmann factors to avoid numerical underflow; this shift is restored in the mean
energy.

With $k_{\mathrm B}=1$, the canonical internal energy and heat capacity are
\begin{align}
  U(T) &= \expval{H}_\beta
  =\frac{1}{Z}\sum_{k,n}m_k E_{n,k}e^{-\beta E_{n,k}},\\
  C_V(T) &= \frac{\partial U}{\partial T}
  =\beta^2\left(\expval{H^2}_\beta-\expval{H}_\beta^2\right).
  \label{eq:capacity}
\end{align}
The temperature scan in Figure~\ref{fig:capacity-T} uses $L=20$, $r=0.2$,
$0.02\leq T\leq2$, and $g=0,-0.1,\ldots,-0.5$. The internal energy rises
smoothly with temperature. Each heat-capacity curve has a broad maximum, whose
position shifts to higher temperature as the deformation becomes stronger. The scan
therefore maps the thermodynamic crossover produced by the same deformation that
moves the scarred spectral weight. At finite size, the broad maximum locates the
temperature range with the strongest energy fluctuations; it is not by itself a
signature of Griffiths behavior or a thermodynamic phase transition.

\begin{figure}[H]
  \centering
  \includegraphics[width=\linewidth]{pxp_capacity_T_L=20_r=0.2.png}
  \caption{Canonical thermodynamics as a function of temperature for $L=20$ and
  $r=0.2$. The left panel shows the total internal energy $U=\expval{H}$; the
  right panel shows the total heat capacity $C_V$. Each curve corresponds to a
  different deformation strength $g$.}
  \label{fig:capacity-T}
\end{figure}

Figure~\ref{fig:capacity-fit} tests the proposed low-temperature form
$C_V(T)=A T^\lambda$. For each $g$, the fit includes temperatures from
$T=0.02$ to one half of the temperature at which $C_V$ reaches its maximum.
The fitted curves do not follow the numerical data closely: the coefficients of
determination range from $R^2=0.845$ to $0.865$, and the fitted exponents depend
strongly on $g$. A finite chain has a discrete low-energy spectrum and generally a
nonzero gap $\Delta$. Its asymptotic low-temperature heat capacity is therefore
activated,
\begin{equation}
  C_V(T)\sim \left(\frac{\Delta}{T}\right)^2 e^{-\Delta/T},
\end{equation}
rather than a power law. At higher temperatures, additional levels contribute as
the system approaches the broad heat-capacity maximum. The selected interval thus
crosses a finite-size crossover instead of isolating a scale-invariant regime. The
reported values of $\lambda$ are window-dependent effective slopes and do not
establish Griffiths scaling.

\begin{figure}[H]
  \centering
  \includegraphics[width=0.88\linewidth]
  {pxp_capacity_fit_L=20_r=0.2.png}
  \caption{Log--log comparison of the canonical heat capacity (solid curves) and
  fits to $C_V=A T^\lambda$ (dashed lines) for $L=20$ and $r=0.2$. Each fit ends
  at one half of the corresponding peak temperature. The legend gives the fitted
  exponent for each deformation strength.}
  \label{fig:capacity-fit}
\end{figure}

Figure~\ref{fig:capacity-g} presents the complementary scan in which temperature is
fixed and $g$ varies from $-0.5$ to $0$. At the lowest temperatures, the internal
energy follows the deformation of the low-energy spectrum while $C_V$ remains close
to zero. At higher temperatures, the heat capacity becomes sensitive to $g$ across
the full interval. Taken together, Figures~\ref{fig:capacity-T} and
\ref{fig:capacity-g} show how the same spectral rearrangement appears along two cuts
through the $(g,T)$ plane. They also establish a baseline for a sharper test proposed
in this work: a Griffiths regime should produce a non-universal low-temperature law
$C_V(T)\sim T^\lambda$ over the same finite interval of $g$ in which slow dynamics
is observed. Testing that statement requires low-temperature log--log fits, several
system sizes, and control of the finite-size gap.

\begin{figure}[H]
  \centering
  \includegraphics[width=\linewidth]{pxp_capacity_g_L=20_r=0.2.png}
  \caption{Canonical thermodynamics as a function of deformation strength for
  $L=20$ and $r=0.2$. The left and right panels show $U(g)$ and $C_V(g)$,
  respectively, at fixed temperatures $T=0.01$, $0.1$, $0.2$, $0.5$, $0.8$,
  and $1$.}
  \label{fig:capacity-g}
\end{figure}

\section{Connected local autocorrelation}

The dynamical calculation starts from a canonical state of the undeformed PXP
Hamiltonian,
\begin{equation}
  \rho_0=\frac{e^{-\beta\Hpxp}}
  {\Tr(e^{-\beta\Hpxp})},
\end{equation}
and evolves operators with the deformed Hamiltonian $H_1=H(g,r)$,
\begin{equation}
  \Xtilde_i(t)=e^{iH_1t}\Xtilde_i e^{-iH_1t}.
\end{equation}
The measured quantity is the site-averaged connected autocorrelation
\begin{equation}
  C_X(t)=\frac{1}{L}\sum_{i=0}^{L-1}\operatorname{Re}\left[
  \Tr\!\left(\rho_0\Xtilde_i(t)\Xtilde_i\right)
  -\Tr\!\left(\rho_0\Xtilde_i(t)\right)
   \Tr\!\left(\rho_0\Xtilde_i\right)
  \right].
  \label{eq:correlation}
\end{equation}
The subtraction removes the disconnected one-point contribution. Because $\rho_0$
need not commute with $H_1$, this protocol is a thermal quench rather than a
stationary equilibrium correlator. The time-dependent one-point function must
therefore be retained in the subtraction. The code evaluates
Eq.~\eqref{eq:correlation} in the eigenbasis of $H_1$, allowing all requested times
to be obtained from precomputed matrix contractions and phase factors. The local
operator $\Xtilde_i$ is a candidate probe of slow constrained dynamics; identifying
the corresponding quasi-good quantum number or effective order parameter remains
part of the Griffiths interpretation.

The current parameter set is $L=16$, $g=-0.4$, $r=0.2$, and $T=1$. Times are
sampled geometrically over $0.1\leq t\leq100$, which provides uniform resolution on
a logarithmic time axis.

\subsection{Power-law fit}

The observed decay is modeled by
\begin{equation}
  C_X(t)=A t^{-\alpha}.
  \label{eq:powerlaw}
\end{equation}
For positive correlation values, Eq.~\eqref{eq:powerlaw} becomes the linear relation
\begin{equation}
  \log C_X(t)=\log A-\alpha\log t.
  \label{eq:logfit}
\end{equation}
Ordinary least-squares regression is applied directly to all correlation samples in
the early-time interval
\begin{equation}
  0.1\leq t\leq t_{\mathrm{fit}},
  \qquad t_{\mathrm{fit}}=50.
\end{equation}
No peak detection or upper-envelope selection enters the fit. The fitted exponent is
minus the regression slope, and the amplitude is the exponential of the intercept.
The plot uses logarithmic axes and marks the fitted time interval with a translucent
shaded region without an edge line. This regression measures whether a power law is
consistent with the selected window. A model test must additionally fit a single
exponential and a stretched exponential to the same untransformed data and compare
their residual structure and window dependence.

Figure~\ref{fig:autocorrelation} shows the correlation data and the fitted power law.
The regression gives $\alpha=0.486$ for this parameter set. Oscillations remain around
the fitted trend, but the dashed line captures the overall decrease across the shaded
time window. Data beyond $t_{\mathrm{fit}}$ are displayed for comparison and do not
enter the regression.

\begin{figure}[H]
  \centering
  \includegraphics[width=0.8\linewidth]
  {pxp_autocorr_avg_L=16_g=-0.4_r=0.2_T=1.0.png}
  \caption{Connected local $X$-autocorrelation for $L=16$, $g=-0.4$,
  $r=0.2$, and $T=1$ on logarithmic axes. Orange points inside the shaded region
  are the data used in the regression over $0.1\leq t\leq50$. The green dashed
  line is the fitted power law $C_X(t)=At^{-\alpha}$ with $\alpha=0.486$; data at
  later times are excluded from the fit.}
  \label{fig:autocorrelation}
\end{figure}

The approximately linear behavior in the log--log representation supports a
power-law description of the finite-size autocorrelation over the selected window.
Such slow relaxation is compatible with Griffiths phenomenology, in which a broad
distribution of relaxation scales replaces a single decay time. The present model is
clean and staggered rather than spatially random, and the fit includes microscopic
early times. The result is therefore an initial diagnostic, not a sufficient
identification of a Griffiths phase. Its physical content must be tested against
alternative decay models, shifted fit windows, additional temperatures and
deformations, and increasing system size.

\section{Tests for a disorder-free Griffiths regime}

The numerical results above suggest a connection between a low-energy scarred
subspace and slow relaxation. The following tests are needed to decide whether that
connection defines a disorder-free Griffiths regime rather than a finite-size
crossover.

\begin{enumerate}
  \item \textbf{Joint dynamical and thermodynamic scaling.}
  The autocorrelation should favor an algebraic decay over a single exponential
  $A e^{-t/\tau}$ and a stretched exponential
  $A\exp[-(t/\tau)^b]$ across a broad intermediate-time window. Independently,
  $C_V(T)$ should exhibit a compatible non-universal power law at low temperature.
  Both behaviors should persist over the same finite interval of $g$ rather than at
  one finely tuned point.

  \item \textbf{Separation between average and typical observables.}
  For a positive observable or response $X$, a broad distribution can be diagnosed
  by comparing
  \begin{equation}
    X_{\mathrm{avg}}=\langle X\rangle,
    \qquad
    X_{\mathrm{typ}}=\exp\!\left(\langle\ln X\rangle\right).
  \end{equation}
  A strong inequality $X_{\mathrm{typ}}\ll X_{\mathrm{avg}}$ would show that the
  mean is dominated by rare contributions. In this clean translation-invariant
  model, a site average cannot reproduce a disorder average. The relevant ensemble
  must first be defined, for example over eigenstates, symmetry sectors, phase-space
  regions, or controlled weak perturbations.

  \item \textbf{Distribution of low-energy gaps.}
  If $\Delta$ samples the relevant low-energy or relaxation gaps, compare its mean
  with
  \begin{equation}
    \Delta_{\mathrm{typ}}
    =\exp\!\left(\langle\ln\Delta\rangle\right).
  \end{equation}
  The hierarchy $\Delta_{\mathrm{typ}}\ll\langle\Delta\rangle$ would indicate a
  broad gap distribution. The ensemble used to construct this distribution must be
  specified before the comparison is meaningful.

  \item \textbf{Negative checks.}
  Failure of a single exponential over a growing time window would distinguish the
  candidate regime from a conventional gapped paramagnet. Persistence of anomalous
  scaling over a finite parameter interval, together with finite-size analysis,
  would distinguish it from an isolated critical point.

  \item \textbf{Microscopic origin of the slow sector.}
  The remaining conceptual task is to identify the approximate quantum number or
  phase-space structure that plays the role normally assigned to a locally ordered
  rare region. Establishing that structure would turn the observed analogy into a
  mechanism.
\end{enumerate}

\section{Summary}

The deformed PXP chain provides a clean setting in which a scarred subspace is shifted
toward low energy. Exact diagonalization identifies eigenstates with anomalously
large $\Ztwo$ overlap, symmetry-resolved spectra determine the canonical energy and
heat capacity, and a thermal quench produces a connected local $X$-autocorrelation
consistent with power-law decay over the fitted time window. These results establish
the finite-size spectral and dynamical ingredients of a possible connection between
quantum many-body scars and Griffiths-like relaxation.

The stronger claim of a disorder-free Griffiths regime remains conditional. It
requires consistent dynamical and heat-capacity exponents over a finite deformation
window, a meaningful separation between average and typical observables, a broad
distribution of gaps, and identification of the approximate structure responsible
for rare slow events. Those tests define the next stage of the study and separate the
proposed mechanism from ordinary crossover behavior or isolated criticality.

\end{document}
